% 01-introduction.tex - Introduction & Credits

\section{Introduction}

\begin{frame}{Beamer vs. PowerPoint}
Compared to PowerPoint, using \LaTeX\ with Beamer is better because:
\begin{itemize}
    \item \textbf{What-You-Mean-Is-What-You-Get}: You write the content, the computer handles the professional typesetting.
    \item \textbf{PDF Portability}: Produces a \texttt{.pdf} directly --- no worries about missing fonts, broken formulas, or mismatched software versions.
    \item \textbf{Consistency}: Extremely easy to maintain consistent layouts, fonts, highlights, and colors.
    \item \textbf{Mathematical Typesetting}: Standard \TeX\ math rendering is the best in class:
    \begin{equation*}
        \mathrm{i}\,\hslash\frac{\partial}{\partial t} \Psi(\mathbf{r},t) =
        -\frac{\hslash^2}{2\,m}\nabla^2\Psi(\mathbf{r},t) + V(\mathbf{r})\Psi(\mathbf{r},t)
    \end{equation*}
\end{itemize}
\end{frame}

\begin{frame}{Template Credits}
This presentation theme has a rich history of contributions:
\begin{itemize}
    \item \textbf{Original SINTEF Theme}: Created by \hrefcol{mailto:federico.zenith@sintef.no}{Federico Zenith}, derived from Håvard Berland's \texttt{beamerthementnu}.
    \item \textbf{PoliTo \& Sapienza Adaptation}: Adapted for Politecnico di Torino by \hrefcol{mailto:mattia.ippoliti@studenti.polito.it}{Mattia Ippoliti} and Sapienza by Andrea Gasparini.
    \item \textbf{College Beamer inspiration}: Liu Qilong's collection is retained as a documented influence; no current file is vendored from that repository.
    \item \textbf{Apograph Library Generalization}: Packaged and generalized by Elia Innocenti (providing robust fallbacks, path cleaning, and dynamic metadata).
\end{itemize}
\vspace{0.3cm}
\begin{center}
    \footnotesize Presentation content carries CC BY 4.0 attribution; the theme code is GPL-3.0-or-later. See \texttt{NOTICE} and \texttt{LICENSES/}.
\end{center}
\end{frame}
